%------------------------------------------------------------------------------%
\section{Integração por substituição simples}
\label{sec:Integracao_substituicao_simples}

%------------------------------------------------------------------------------%
\begin{teorema}{A Regra da substituição}{teo:substituicao}
  Se $u=g(x)$ for uma função derivável cuja imagem é um intervalo
  $I$ e $f$ for contínua em $I$, então
  \[
    \int f(g(x)) g'(x) dx= \int f(u) du
  \]
\end{teorema}
%------------------------------------------------------------------------------%

%------------------------------------------------------------------------------%
\begin{example}
  Encontre \(\int 4x^3 \cos(x^4+2)dx\).
  \solution
  Note que, podemos escrever a integral
  \[
    \int 4x^3 \cos(x^4+2)dx =  \int  \cos(x^4+2) 4x^3dx
  \]
  como o produto de duas funções $f(x)=cos(x^4+2)$ e $h(x)= 4x^3$. Note que $f$
  é uma composição de duas funções e a função $h$ é exatamente a derivada da
  função ``de dentro'' da composição, ou seja, a deribada de $x^4+2$. Portanto,
  nesse caso, a regra da substituição simples parece se encaixar bem.

  Fazendo a substituição $u=x^4+2$ temos que a sua diferencial é $du=4x^3dx$ e
  usando a regra da substituição temos
  \begin{align*}
    \int 4x^3 \cos(x^4+2)dx
    & =  \int  \cos(x^4+2) 4x^3dx= \int cos(u) du \\
    & = \sin(u)+C \\
    & = -\sin(x^4+2)+C
  \end{align*}
\end{example}
%------------------------------------------------------------------------------%

\begin{itemize}
  \item Note que, quando usamos a regra da substituição, após integrarmos com
      relação a variável $u$ temos que voltar para a variável original.

  \item Quando usamos a regra da substituição nossa intenção é transformar uma
      integral complicada em uma mais simples de ser calculada.

\end{itemize}

%------------------------------------------------------------------------------%
\begin{example}
  Calcule \(\int \frac{x}{\sqrt{1-4x^2}}dx\)
  \solution
  Note que se fizermos a substituição $u=1-4x^2$, então
  $du=-8x dx \Leftrightarrow xdx = -\frac{1}{8} du$.
  Assim
  \begin{align*}
    \int \frac{x}{\sqrt{1-4x^2}}dx
    & = -\frac{1}{8} \int \frac{1}{\sqrt{u}}du \\
    & = -\frac{1}{8} \int u^{-\frac{1}{2}} du \\
    & = -\frac{1}{8}(2 \sqrt{u})+C \\
    & = -\frac{1}{4}\sqrt{1-4x^2}+C
  \end{align*}
\end{example}
%------------------------------------------------------------------------------%

%------------------------------------------------------------------------------%
\begin{example}
  Calcule \(\int \tan(x) dx\)
  \solution
  Sabemos que
  \[
    \tan(x)=\frac{\sin(x)}{\cos(x)},
  \]
  portanto
  \[
    \int \tan(x) dx= \int \frac{\sin(x)}{\cos(x)} dx
  \]
  Isso nos sugere fazer a substituição $u=\cos(x)$, com isso,
  $du=-\sin(x)dx \Leftrightarrow -du= \sin(x)dx$.
  Assim, usando a regra da substituição simples
  \begin{align*}
    \int \tan(x) dx
    & = \int \frac{\sin(x)}{\cos(x)} dx \\
    & = \int \frac{1}{u} -du \\
    & = - \int \frac{1}{u} du \\
    & = - \ln \abs{\cos(x)}+C
  \end{align*}
\end{example}
%------------------------------------------------------------------------------%

%------------------------------------------------------------------------------%
\begin{example}
  Encontre \(\int_1^e \frac{\ln(x)}{x}dx\)
  \solution
  Vamos resolver esse exercício de duas maneiras.

  \emph{1\textsuperscript{o} modo}:
  Primeiro vamos calcular a integral indefinida
  $\int \frac{\ln(x)}{x} dx$ e depois usaremos o
  Teorema Fundamental do Cálculo parte 2 para
  calcular a integral indefinida.

  Para calcular a integral indefinida iremos usar a substituição
  $u=\ln(x)$ e com isso, $du= \frac{1}{x}dx$.
  \begin{align}
    \int \frac{\ln(x)}{x} dx
    & = \ln(x) \frac{1}{x} dx \\
    & = \int u du \\
    & = \frac{u^2}{2}+C \\
    & = \frac{(\ln(x))^2}{2}+C
  \end{align}

  Usando o Teorema fundamental do Cálculo parte 2:
  \begin{align*}
    \int_1^e \frac{\ln(x)}{x} dx
    & = \left[ \frac{(\ln(x))^2}{2} \right]_1^e    \\
    & = \frac{(\ln(e))^2}{2} -\frac{(\ln(1))^2}{2} \\
    & = \frac{1}{2}
  \end{align*}
\end{example}
%------------------------------------------------------------------------------%

Algumas observações importantes
\begin{itemize}
  \item Seja $F$ a primitiva mais geral de $f$, ou seja,
  \[
    \int f(x)dx= F(x)+c,
  \]
    então
  \[
    \int f(kx) dx= \frac{F(kx)}{k}+C.
  \]
  \item Além disso, \(\int f(kx+b)dx= \frac{F(kx+b)}{k}+C\).
  \item \(\int \cos(5x) dx= \frac{\sin(5x)}{5}+C\).
  \item \(\int e^{8x}= \frac{e^{8x}}{8}+C\).
  \item \(\int \frac{1}{x+3}dx= \ln\abs{x+3}+C\).
\end{itemize}
Essas observações podem ser demonstradas usando substituição simples.

%------------------------------------------------------------------------------%
